The three defects discussed in this report are violations of standard algebraic laws from category theory. This section fixes notation and restates the relevant definitions precisely enough to make the violations, and the proofs that the fixes restore the laws, unambiguous. The definitions follow Barr and Wells~\cite{BarrWells1990}; the monadic treatment of \texttt{Option}/\texttt{Either} follows Moggi~\cite{Moggi1991} and Wadler~\cite{Wadler1992}.

\subsection{Categories, Functors, and Natural Transformations}

A \emph{category} $\mathcal{C}$ consists of a collection of objects $\mathrm{Ob}(\mathcal{C})$, and for every pair of objects $u,v$ a set of morphisms $\mathrm{Hom}(u,v)$, together with an associative composition operator $\circ$ and, for every object $u$, an identity morphism $\mathrm{id}_u \in \mathrm{Hom}(u,u)$ satisfying
\[
  \mathrm{id}_v \circ f = f = f \circ \mathrm{id}_u \qquad \text{for all } f \in \mathrm{Hom}(u,v),
\]
\[
  (h \circ g) \circ f = h \circ (g \circ f) \qquad \text{whenever the composites are defined.}
\]
In \texttt{CategoryTheory.php}, the directed travel graph $G=(V,E)$ induces exactly such a category $F(G)$ (Section~\ref{sec:bug3} develops this in detail): $\mathrm{Ob}(F(G))=V$, and $\mathrm{Hom}(u,v)$ is the set of finite, directed walks from $u$ to $v$ along edges of $E$, with composition given by path concatenation and $\mathrm{id}_u$ the empty walk at $u$. This is the \emph{free category} generated by $G$: it has no morphisms beyond those forced to exist by the graph and the two axioms above.

A \emph{functor} $F : \mathcal{C} \to \mathcal{D}$ maps every object $u$ of $\mathcal{C}$ to an object $F(u)$ of $\mathcal{D}$ and every morphism $f : u \to v$ to a morphism $F(f) : F(u) \to F(v)$, subject to
\[
  F(\mathrm{id}_u) = \mathrm{id}_{F(u)}, \qquad F(g \circ f) = F(g) \circ F(f).
\]
\texttt{RouteList} implements this directly for the category of PHP values under function application: \texttt{RouteList::fmap(\$f)} maps every element with \texttt{\$f}, and \texttt{category\_laws\_test.php} checks both equations against the \texttt{discount}/\texttt{tax} transformations used throughout the booking pipeline.

A \emph{natural transformation} $\eta : F \Rightarrow G$ between two functors $F,G:\mathcal{C}\to\mathcal{D}$ is a family of morphisms $\eta_u : F(u) \to G(u)$, one per object $u$ of $\mathcal{C}$, such that for every $f : u \to v$ the naturality square commutes:
\[
  G(f) \circ \eta_u = \eta_v \circ F(f).
\]
\texttt{RouteList::headOption} is such a transformation from the \texttt{RouteList} functor to the \texttt{Option} functor: picking the head of a list and then mapping a function over the resulting \texttt{Option} must agree with mapping the function over the whole list first and then picking the head (both results are checked for equality in the suite, including the empty-list edge case where both sides must collapse to \texttt{None}).

\subsection{Monads, the Kleisli Category, and \texttt{Option}/\texttt{Either}}

A \emph{monad} on $\mathcal{C}$ is a functor $M$ together with two operations, $\mathrm{pure} : a \to M(a)$ and $\mathrm{bind} : M(a) \to (a \to M(b)) \to M(b)$ (written infix as $\mathbin{>\!\!>\!\!=}$), satisfying three laws for all $a$, $m : M(a)$ and Kleisli arrows $f,g$:
\[
  \text{(Left identity)}\quad \mathrm{pure}(a) \mathbin{>\!\!>\!\!=} f \;=\; f(a),
\]
\[
  \text{(Right identity)}\quad m \mathbin{>\!\!>\!\!=} \mathrm{pure} \;=\; m,
\]
\[
  \text{(Associativity)}\quad (m \mathbin{>\!\!>\!\!=} f) \mathbin{>\!\!>\!\!=} g \;=\; m \mathbin{>\!\!>\!\!=} \big(\lambda x.\, f(x) \mathbin{>\!\!>\!\!=} g\big).
\]
These three laws are exactly what make Kleisli arrows $a \to M(b)$ composable into a category (the \emph{Kleisli category} of $M$), with $\mathrm{pure}$ as identity and $\mathbin{>\!\!>\!\!=}$-chaining as composition. \texttt{Option} (the \texttt{Some}/\texttt{None} pair, PHP's analogue of \texttt{Maybe}) and \texttt{Either} (\texttt{Right}/\texttt{Left}, used for the short-circuiting validation pipeline in \texttt{index.php}) are both implemented as such monads. \textbf{The right-identity law is precisely the law violated by the first defect}: if $\mathrm{pure}$ does not act as a faithful identity on every value of the carrier type -- in particular on \texttt{bool} -- then $m \mathbin{>\!\!>\!\!=} \mathrm{pure}$ can produce a different \texttt{Option} than $m$ itself, and the Kleisli category stops being a category at all.

\subsection{Monoids and Strict Monoidal Functors}

A \emph{monoid} $(M, \oplus, e)$ is a set $M$ with an associative binary operation $\oplus$ and an identity element $e$:
\[
  (a \oplus b) \oplus c = a \oplus (b \oplus c), \qquad e \oplus a = a = a \oplus e.
\]
\texttt{TripMetrics} is declared a \texttt{Monoid} over $(\text{cost}, \text{duration}, \text{points})$ with \texttt{TripMetrics::empty()} as $e$ and \texttt{append()} as $\oplus$; it is used to fold the per-step costs of a booked itinerary into a single total. \textbf{The associativity law is precisely the law violated by the second defect}: Section~\ref{sec:bug2} exhibits concrete values for which $(a\oplus b)\oplus c \neq a\oplus(b\oplus c)$ under the original implementation.

A monoid homomorphism $M : (\mathcal{C}, \circ) \to (M,\oplus)$ from a category's morphisms into a monoid, satisfying $M(g \circ f) = M(f) \oplus M(g)$ and $M(\mathrm{id}) = e$, is called a \emph{strict monoidal functor}. \texttt{FreeCategoryPath::evaluateMetrics()} is exactly this: it evaluates a composed path in $F(G)$ by folding the \texttt{TripMetrics} monoid over the path's edges, and \texttt{category\_laws\_test.php} checks that evaluating a composed path agrees with evaluating and then combining its parts -- a check that is only meaningful if the underlying monoid actually associates, which is why the fix to \texttt{TripMetrics} in Section~\ref{sec:bug2} is a precondition for this check to mean anything at all.

\subsection{Error Propagation as a Categorical Obligation}

Short-circuiting a computation on failure is not merely a convenient control-flow idiom here: it is the \texttt{Either} monad's \texttt{bind} doing exactly what the monad laws require of it. If a step in a Kleisli chain produces \texttt{Left}, every subsequent \texttt{bind} in the chain must propagate that \texttt{Left} unchanged, by the definition of \texttt{Either::bind}. \textbf{The third defect is, in this sense, not a violation of one of the three named laws above, but a violation of the obligation those laws impose}: \texttt{FreeCategoryPath::compose} correctly computes a \texttt{Left} when two paths do not form a connected walk in $F(G)$, but \texttt{Itinerary::advance()} -- a plain method, not itself a Kleisli arrow -- read that \texttt{Either} and then discarded the failure instead of propagating it, breaking the chain of custody that the rest of the pipeline relies on. Section~\ref{sec:bug3} develops this precisely.
